%% LyX 2.3.3 created this file.  For more info, see http://www.lyx.org/.
%% Do not edit unless you really know what you are doing.
\documentclass{article}
\usepackage[utf8]{inputenc}
\usepackage{geometry}
\usepackage{longtable}
\geometry{verbose,tmargin=1in,bmargin=1in,lmargin=1in,rmargin=1in}
\setlength{\parskip}{\medskipamount}
\setlength{\parindent}{0pt}
\usepackage{amsmath}
\usepackage{url}
\usepackage{amssymb}
\usepackage{setspace}
\onehalfspacing
\begin{document}
\begin{center}
{\Large{}Problem set \#2}{\Large\par}
\par\end{center}

\begin{center}
Due on October 21st. Send solution to econ628vse@gmail.com \\
Answers must be typed. You can work in teams under the constraint: Max(Team Size)=2 \\

\par\end{center}

\bigskip{}
\begin{enumerate}

  \item \textbf{Weighted Least Squares and OverID}
Download the CPS extract and the build do-file from canvas.

a) Regress lnwage on
dummies for the categories of the variables agecat, educcat, and sex. Use main effects only -- i.e. no interactions. 

b) Construct a
dataset of means and number of observations of the variable lnwage
within cells defined by combinations of the variables agecat, educcat,
and sex. To keep things simple don't use the person weights.

c) Run a weighted least squares (WLS) regression of cell wages on dummies for age, education,
and sex (main effects only) using the number of observations in each
cell as a weight. How do the coefficients compare to what you got
in a)?

d) If wages are $iid$ within cell, what is the variance of each cell
mean?

e) Construct an estimate of the variance of each cell
mean. Do not assume the data are homoscedastic across cells.

f) Run a WLS regression of cell wages on dummies for age, education,
and sex (main effects only) using the inverse of your estimate of
the cell variance as a weight.

g) Compute a over-id specification test (as described in class) from the weighted sum of squared residuals
of this regression. What do you conclude about the fit of the main effects model?

h) Run a WLS regression of cell wages on main effects for age, education,
and sex plus all two-way interactions (e.g. age$\times$education,
education$\times$sex, and age$\times$sex).

i) Compute the over-id specification test. Can you reject this model at the
5\% level?

j) Compute the predicted cell means from the model you estimated in
h). Graph the predicted and actual cell means against age by subgroup.
Look for signs of misspecification.

k) Can you find a simple regression specification that passes the
over-id test?

\item \textbf{MLE:} Let $Y_{i}$ be a dichotomous $0/1$ random variable and
$X_{i}$ a $K\times1$ vector of predictor variables. Consider the
Logit log-likelihood:

\begin{eqnarray*}
l_{i} & = & Y_{i}\ln\Lambda\left(X_{i}^{\prime}\beta\right)+\left(1-Y_{i}\right)\ln\left[1-\Lambda\left(X_{i}^{\prime}\beta\right)\right]\\
\Lambda\left(X_{i}^{\prime}\beta\right) & = & \frac{\exp\left(X_{i}^{\prime}\beta\right)}{1+\exp\left(X_{i}^{\prime}\beta\right)}
\end{eqnarray*}

Denote the maximizer of $E[l_{i}]$ as $\beta_{ML}$ and the maximizer
of $\frac{1}{N}\sum\limits _{i}l_{i}$ as $\widehat{\beta}_{ML}$.

a) Derive the score of the logit log-likelihood

b) Using iterated expectations provide a discussion of the moment
conditions identifying $\beta_{ML}$.

c) Define 
\[
\beta_{NLLS}=\underset{\beta}{\arg\min}E\left[\left(Y-\Lambda\left(X_{i}^{\prime}\beta\right)\right)^{2}\right]
\]

Derive the population moment conditions identifying $\beta_{NLLS}$

d) Under what conditions will $\beta_{NLLS}$ and $\beta_{ML}$ coincide?

e) Derive an expression for the asymptotic variance of $\widehat{\beta}_{ML}$
under the assumption of proper model specification

f) Derive an expression for the asymptotic variance without the assumption
of proper specification but with the assumption of independence across
observations.

    \item \textbf{SML:} Consider the following random coefficient binary choice
model:

\[
Y_{i}=1\left[b_{i}X_{i}+\varepsilon_{i}>0\right]
\]

where $\varepsilon_{i}\sim Logistic(0,1)$ and $b_{i}\sim N\left(\mu,\sigma^{2}\right)$.

a) Write an expression for the simulated log likelihood and its derivatives
with respect to $\mu$ and $\sigma$.

b) Load the dataset ``logit". The first column of the matrix
``data'' gives observations on $Y_{i}$ and second on $X_{i}$.
Generate a 2000x1000
matrix $V$ of draws from a standard  normal distribution. The columns of this matrix index simulation draws while the rows index
observations.

c) Write a program ``SML" that takes as arguments the data vectors,
the matrix $V$ of random variates, and the parameters $\left(\mu,\ln\sigma\right)$
and returns the negative of the simulated log likelihood and its analytic
gradient (or score). 


d) Maximize the simulated likelihood and
 compute the Hessian (second derivative of the objective function with respect to $\mu$ and $\sigma$) of the SML criterion. Make sure that your maximization routine uses the analytical gradient that your provided in the previous part.
Report your SML parameter estimates of $\left(\mu,\ln\sigma\right)$.
Verify that your estimates are not sensitive to starting values.

e) Report standard errors for $\left(\mu,\ln\sigma\right)$ computed
under the assumption of proper specification.

f) Report robust standard errors based upon the sandwich matrix.

 
\item \textbf{Cluster Robust Standard Errors} Generate 100 clusters and 100 observations
within each cluster from the following DGP: 
\begin{eqnarray*}
Y_{ic} & = & \nu_{c}+D_{ic}\eta_{c}+\varepsilon_{ic}\\
D_{ic} & = & I\left[D_{ic}^{\ast}>.5\right]\\
D_{ic}^{\ast} & \sim & U\left[0,1\right]\\
\eta_{c} & \sim & N\left(1,\sigma_{\eta}^{2}\right),\nu_{c}\sim N\left(0,1\right)\\
\varepsilon_{ic} & \sim & N\left(0,1\right)
\end{eqnarray*}
This DGP can be thought of as representing the effect of some randomized
treated $D_{ic}$ assigned at the individual level with probability
$\frac{1}{2}$ which has heterogeneous effects $\eta_{c}$ across
clusters (sites) but an average effect of 1. There is also site heterogeneity
in untreated outcomes as captured by the $\nu_{c}$.

a) Simulate data from this DGP setting $\sigma_{\eta}^{2}=1$ and fixing the seed simulator to 1234 for replication purposes. Regress
$Y_{ic}$ on $D_{ic}$. What is the standard error of the
estimated coefficient on $D_{ic}$ under homoskedasticity? What is the ``robust\textquotedblright{} or White
standard error? What is the ``clustered\textquotedblright{} standard
error?


b) Repeat step a) \underline{using the same simulated values} that you used in part (a) but this time setting $\sigma_{\eta}^{2}=0$.
What is the ``robust\textquotedblright{} standard error? What is
the ``clustered\textquotedblright{} standard error?

c) Comment on the differences between your answers in a) and b).

d) Now perform a Monte
Carlo exercise generating data from the above DGP 1000 times when
$\sigma_{\eta}^{2}=1$. Each time perform a t-test of the (true) null
that the population regression coefficient on $D_{c}$ equals one
based upon: i) the robust standard error from the microlevel regression,
ii) the clustered standard error (iii) a block-level bootstrap where the ``block" is going to be given by a given cluster. Use a significance level of 5\% and
in each simulation record whether the test in question rejects. What
fraction of the time does each test reject the null that the regression
coefficient is equal to 1?

e) Now conduct a Monte Carlo where the $\left\{ \eta_{c}\right\} _{c=1}^{100}$
are fixed in repeated draws as would be the case if we considered
re-running the experiment on the same 100 sites in our study. What
fraction of the time does each test reject the true null that the
regression coefficient equals $\frac{1}{100}\sum_{c=1}^{100}\eta_{c}$? 

f) Compare and contrast your answers from parts d) and e). What does
this say about the appropriate level of clustering in randomized experiments?

\item (\textbf{Bootstrap OLS)} Download the Stata dataset ``bootstrap.dta''
from the course website. This dataset has a binary outcome y, a binary
regressor of interest D, a cluster level control X, and a micro control
X2. There are 20 clusters. Set the seed to ``123''. Use 1,000 reps
for all bootstrap procedures.

a) Regress y on D, X, and X2 clustering by the variable ``id''.
Report the clustered standard error for the coefficient on D.

b) What is the p-value for the null hypothesis that the coefficient
on D equals zero?

c) Block bootstrap the regression (make sure to use the ``, cluster(id)''
switch). What is the bootstrap standard error on the coefficient on
D? Use this standard error to compute a new p-value.

d) Block bootstrap the clustered t-statistic from the regression (make
sure to use both the cluster() switch and the idcluster() switch).
Use a symmetric bootstrap percentile t-test to compute a new p-value.

e) Use the Wild bootstrap procedure of Cameron, Gelbach, and Miller
(2008) to test the null that the coefficient on D equals zero. Use
Rademacher weights and make sure to ``impose the null'' hypothesis
on the coefficient on D in the manner they describe. Report the corresponding
p-value.

f) Discuss the different answers given by parts b)-e) in light of
econometric theory. Which answer do you think is most accurate?

\item  (\textbf{Bootstrap Probit}) 

a) Using bootstrap.dta again, fit a Probit of y on D, X, and X2. Report
the clustered standard error for the coefficient on D. What is the
p-value for the hypothesis that the coefficient on D equals zero?

b) Block bootstrap the clustered t-statistic from the Probit. Use
a symmetric bootstrap percentile-t test to compute a new p-value.

c) Conduct a cluster robust score test of the null hypothesis that
the coefficient on D equals zero. Report the p-value.

d) Now block bootstrap the clustered t-statistic.
Use a symmetric bootstrap percentile t-test to compute a new p-value.

e) Which of these p-values do you think is most accurate? Discuss
the tradeoff involved in choosing between alternative bootstrap methods.

\item (\textbf{Reweighting}) Go to David Autor's webpage and download
the cleaned 1979 and 1997 MORG files from:

https://economics.mit.edu/people/faculty/david-h-autor/data-archive

the paper is Trends in U.S. Wage Inequality: Revising the Revisionists.


a) Read through the cleaning programs associated with the MORG files.
Describe the important decisions being made.

b) Make a new variable ``lnhr\_wage'' which is the log of the ``hr\_wage''
variable. Restrict attention to the cleaned sample with ``hr\_wage\_sample==1''.
Use the ``kdensity'' command to plot kernel density estimates of
lnhr\_wage using the variable wgt\_hrs (which is the product of the
sample weights and the number of hours worked) as an ``aweight''.
Make density plots for both years on the same grid.

c) How has the wage distribution changed over time?

d) A number of demographic shifts occured between 1979 and 1997 that
might alter the aggregate wage distribution even if the distribution
for each demographic subgroup remained unaltered. Repeat the exercise
in b) separately for black men age 25-50 and for white women age 25-50.
How are the patterns different?

e) One way of adjusting the distribution for compositional differences
is to use propensity score reweighting.

Append the 1979 and 1997 files together and create a dummy variable
equal to one for observations in the 1997 sample. Estimate a logit
of the 1997 dummy on a dummy for female, a black dummy, an ``other
race'' dummy, and a quadratic polynomial in age. Use the ``aweights''
when estimating the logit.

f) Use the propensity score implied by this logit to reweight the
1997 wage distribution so that it has the 1979 demographic distribution.
Make a kernel density plot of the actual 1979 and 1997 wage distributions
against the ``counterfactual'' reweighted 1997 distribution.

g) Use the propensity score to reweight the 1979 wage distribution
so that it has the 1997 demographic distribution (if you get stuck
read the Dinardo, Fortin, and Lemieux paper). Make a kernel density
plot of the actual 1979 and 1997 wage distributions against the ``counterfactual''
reweighted 1979 distribution. Be sure to use the product of the wgt\_hrs
variable and the propensity score based weights when computing the
kernel density.

h) Based upon your analysis what would you say was the effect of the
demographic changes between 1979 and 1997 on the evolution of the
aggregate wage distribution?

i) Examine the suitability of your propensity score specification
by comparing the reweighted mean, standard deviation, and the 10th,
25th, 50th, 75th, and 90th percentiles of age in the 1997 sample to
the equivalent moments in the 1979 sample.

j) Do the same exercise restricting attention to blacks.

k) Try experimenting with different specifications for age in the
propensity score logit and checking how the results of i) and j) change.

l) What do you conclude about the suitability of the propensity score
specification?

\item (\textbf{Duflo, Hanna and Ryan (2012)}) This question asks you to solve and estimate a simplified version of the model from Duflo, Hanna and Ryan (2012). Consider a teacher deciding between work and leisure in two time periods, \(t \in \{0,1\}\). Let \(L_t \in \{0,1\}\) indicate whether the teacher takes leisure in period \(t\). The teacher's flow utility in period \(t\) is given by

\[
U_t = \beta C_t + L_t \cdot (\mu + \varepsilon_t),
\]

where \(C_t\) is consumption in period \(t\). All income and consumption is assumed to happen in period 2, so \(C_1 = 0\). In period 2, the teacher receives a salary payment of 5, and an extra bonus of 5 if she worked in both periods (that is, if \(L_1 = L_2 = 0\)). Assume that the teacher has a discount factor \(\delta \in [0,1]\).

\begin{enumerate}
    \item What are the state variables in period 2? Write an expression for the teacher's value function in period 2.
    
    \item Suppose that \(\varepsilon_t \sim \mathcal{N}(0,1)\). Compute the teacher's expected second-period payoff. Use this to write an expression for the teacher’s value function in period 1.
    
    \item Suppose you have a sample of teachers' choices \((L_{1i}, L_{2i})\) for \(i \in \{1, \dots, N\}\). Write an expression for the likelihood of teacher \(i\)'s choice sequence, \(\Pr[(L_1 = L_{1i}) \cap (L_2 = L_{2i})]\), as a function of the parameters \(\beta\), \(\mu\) and \(\delta\). Use this to write an expression for the sample log-likelihood.
    
    \item Download the data set \texttt{``dhr\_example.csv"} from the course website. This data set includes 800 observations on teachers' choice sequences \((L_{1i}, L_{2i})\). Import the data into a statistical software package. Estimate the parameters \((\beta, \mu, \delta)\) by maximum likelihood, and report your parameter estimates.
    
    \item Estimate an alternative version of the model where you assume that teachers are not forward-looking. Use your results from the two models to test the null hypothesis that teachers are not forward-looking. Discuss intuitively the features of the data that allow you to determine whether teachers are forward-looking.

    \item Estimate an alternative version of the model where $\mu\sim N(\mu_{1};\sigma_{1})$. How does this version of the model differs from the previous one? Write down the associated likelihood and provide estimate of $(\beta,\delta, \mu_{1},\sigma_{1},$) along with their standard errors using SML. Make sure to explain how you derived the standard errors for your coefficients. 

\end{enumerate}

\end{enumerate}

\end{document}
